\chapter{配置与选项}
\chapterlead{\texttt{configuration.nix} 是一份声明：系统应当是什么样子。本章说明这份文件的结构、如何查到可用选项、模块如何组合，以及 \texttt{stateVersion} 为什么在安装后不再改动。}

\section{起步：查一个选项}

系统上所有可配置项都由模块声明，可以用 \texttt{nixos-option} 查询。

\begin{shellsession}
$ nixos-option services.openssh.enable
$ nixos-option services.openssh.settings.PermitRootLogin
$ nixos-option users.users
\end{shellsession}

输出包含三项：当前值、默认值、以及该选项的说明与类型。找不到某个选项时，先确认名字拼写，再用下面的方式搜索。

\begin{shellsession}
$ nix search nixpkgs openssh
$ man configuration.nix    # 全部选项的离线手册，可搜索
\end{shellsession}

\section{文件结构}

一份最小的 \texttt{configuration.nix} 是一个函数，接受一组参数，返回一个属性集。

\begin{nixcode}
{ config, pkgs, lib, ... }:

{
  imports = [ ./hardware-configuration.nix ];

  boot.loader.systemd-boot.enable = true;
  networking.hostName = "nixos-lab";
  time.timeZone = "Asia/Shanghai";

  environment.systemPackages = with pkgs; [ git vim ];

  system.stateVersion = "26.11";
}
\end{nixcode}

四点需要注意。花括号外的参数列表是模块收到的上下文，\texttt{config} 可读取其他选项的值，\texttt{pkgs} 是软件包集合，\texttt{...} 表示接受未列出的参数。\texttt{imports} 引入其他模块，最后合并成一个配置。每一项以分号结束，嵌套层级用点号表示。\texttt{with pkgs;} 把包集合的名字引入当前作用域，省去反复写 \texttt{pkgs.}。

\section{选项的类型}

选项有类型，赋值的形状必须匹配。

\begin{booktable}{\textwidth}{J{0.846} J{1.154}}
\toprule
类型 & 赋值形式 \\
\midrule
布尔 & \texttt{true} 或 \texttt{false} \\
字符串 & \texttt{"nixos-lab"} \\
整数 & \texttt{22} \\
列表 & \texttt{[ "wheel" "networkmanager" ]} \\
属性集 & \texttt{\{ ... \}} \\
路径 & \texttt{./hardware-\allowbreak configuration.nix} \\
包 & \texttt{pkgs.git} \\
\bottomrule
\end{booktable}

类型不匹配时，构建在求值阶段就失败，报错会指出选项名与期望类型。这类错误不需要进入编译，改对即可。

\section{合并规则}

多个模块给同一个选项赋值时，结果由选项的合并函数决定。列表默认相加，属性集默认递归合并，单值默认冲突。

\begin{nixcode}
# 模块 A
environment.systemPackages = [ pkgs.git ];

# 模块 B
environment.systemPackages = [ pkgs.vim ];
# 结果：[ pkgs.git pkgs.vim ]
\end{nixcode}

单值冲突时构建失败并指出两处来源，这是有意的设计：两种说法不可能同时成立，让人来选择，而不是悄悄覆盖。

需要强制覆盖时用 \texttt{lib.mkForce}，需要取消某项时用 \texttt{lib.mkDefault} 降低优先级，或者在列表里用 \texttt{lib.mkBefore} 调整顺序。这些函数应当少用：它们让配置的最终值难以推断。

\section{用 lib 写条件}

\texttt{lib} 提供了一批小函数，最常用的是条件与生成。

\begin{nixcode}
{ config, pkgs, lib, ... }:

{
  # 只有笔记本才启用的配置
  services.thermald.enable = lib.mkIf (config.networking.hostName == "laptop") true;

  # 由列表生成多个用户的公钥
  users.users.deploy.openssh.authorizedKeys.keys =
    map (k: "ssh-ed25519 ${k} deploy") [
      "AAAA...1"
      "AAAA...2"
    ];
}
\end{nixcode}

\texttt{lib.mkIf 条件 值} 在条件为假时把这一项从配置中移除，因此依赖它的其他选项也不会生效。这个模式在同时维护台式机与笔记本的配置时很有用。

\section{stateVersion}

\texttt{system.stateVersion} 记录这份配置初次安装时的版本，例如 \texttt{"26.11"}。它影响少数无法自动迁移的数据格式，比如数据库目录结构与用户密码哈希方案。

装好之后不要改这个值，即使升级到了新版本。它表达的是“我的数据是按哪个版本初始化的”，不是“我现在的版本”。改动它可能让已有的数据被按新格式解读，导致服务起不来。新装的机器写当前版本即可。

\section{把配置拆成多个文件}

配置增长后，按主题拆分比堆在一个文件里好维护。

\begin{nixcode}
# configuration.nix
{
  imports = [
    ./hardware-configuration.nix
    ./modules/desktop.nix
    ./modules/dev.nix
  ];
}
\end{nixcode}

被引入的文件与主文件结构相同，都是接受参数、返回属性集的函数。拆分的原则是：一个文件负责一类事，且能单独删除而不影响其他部分。

\section{练习}

\begin{enumerate}
  \item 用 \texttt{nixos-option} 查出 \texttt{services.openssh} 下三个子选项的当前值与默认值。
  \item 把 \texttt{environment.systemPackages} 拆到 \texttt{modules/packages.nix}，重建后确认包仍然存在。
  \item 用一个 \texttt{lib.mkIf} 让某项配置只在主机名为 \texttt{nixos-lab} 时生效，改主机名后重建，观察差异。
  \item 故意给布尔选项赋一个字符串，阅读报错，说明它指出了哪些信息。
  \item 在 \texttt{man configuration.nix} 中搜索 \texttt{firewall}，找出开放端口的选项名。
\end{enumerate}

\begin{checkpoint}
本章结束时，应能：
\begin{itemize}
  \item 用 \texttt{nixos-option} 与离线手册查到任意选项；
  \item 说明模块的参数列表、\texttt{imports} 与属性集的作用；
  \item 判断七种常见类型的赋值形式；
  \item 解释列表与属性集的合并规则，知道 \texttt{mkForce} 的代价；
  \item 说明 \texttt{stateVersion} 的含义，并知道为何不能改。
\end{itemize}
\end{checkpoint}
